% ============================================================
%  C: Como Programar -- 6ª Edição | Deitel & Deitel
%  Capítulo 2 -- Introdução à Programação em C
%  FAZENDO A DIFERENÇA (2.32 e 2.33)
%  Abordagem incremental: Caps. 1 + 2
% ============================================================
\documentclass[12pt,a4paper]{article}
\usepackage[utf8]{inputenc}
\usepackage[T1]{fontenc}
\usepackage[brazil]{babel}
\usepackage{geometry}
\geometry{top=2.5cm,bottom=2.5cm,left=3cm,right=3cm}
\usepackage{parskip}\setlength{\parskip}{6pt}\setlength{\parindent}{0pt}
\usepackage{xcolor,tcolorbox,enumitem,listings,fancyhdr,titlesec,hyperref,booktabs,array,amsmath}
\tcbuselibrary{skins,breakable}

\definecolor{deitelblue}{RGB}{0,70,127}
\definecolor{deitelgray}{RGB}{240,242,245}
\definecolor{deitelgreen}{RGB}{0,110,60}
\definecolor{deitellightblue}{RGB}{220,235,250}
\definecolor{deitelred}{RGB}{160,20,20}
\definecolor{deitellorange}{RGB}{200,90,0}
\definecolor{saude}{RGB}{150,0,60}
\definecolor{saudelightred}{RGB}{245,215,225}
\definecolor{ecogreen}{RGB}{0,120,50}
\definecolor{ecolightgreen}{RGB}{210,240,220}

\newtcolorbox{boaprática}[1][]{colback=deitellightblue,colframe=deitelblue,
  fonttitle=\bfseries\small,title={Boa prática de programação},breakable,#1}
\newtcolorbox{erroco}[1][]{colback=red!5,colframe=deitelred,
  fonttitle=\bfseries\small,title={Erro comum de programação},breakable,#1}
\newtcolorbox{respostabox}[1][]{colback=deitelgray,colframe=deitelblue!60,
  fonttitle=\bfseries\small\color{deitelblue},title={Resposta detalhada},breakable,#1}
\newtcolorbox{enunciadoBox}[1][]{colback=white,colframe=deitelblue,
  fonttitle=\bfseries,title={#1},breakable}
\newtcolorbox{saida}[1][]{colback=black!88,colframe=deitelblue,
  fontupper=\ttfamily\small\color{green!70},title={\small\color{white}Saída do programa},breakable,#1}
\newtcolorbox{atencao}[1][]{colback=yellow!10,colframe=deitellorange,
  fonttitle=\bfseries\small,title={Limitação do Capítulo~2},breakable,#1}
\newtcolorbox{pesquisa}[1][]{colback=ecolightgreen,colframe=ecogreen,
  fonttitle=\bfseries\small,title={Pesquisa externa realizada},breakable,#1}

\setlist[enumerate,1]{label=\textbf{\alph*)},leftmargin=2em}
\setlist[itemize]{leftmargin=2em}

\lstset{
  basicstyle=\ttfamily\small,backgroundcolor=\color{deitelgray},
  frame=single,framerule=0.4pt,rulecolor=\color{deitelblue!40},
  breaklines=true,showstringspaces=false,
  keywordstyle=\color{deitelblue}\bfseries,
  commentstyle=\color{deitelgreen}\itshape,
  stringstyle=\color{deitelred},
  language=C,numbers=left,numberstyle=\tiny\color{gray},
  xleftmargin=18pt,xrightmargin=5pt,stepnumber=1,
}

\pagestyle{fancy}\fancyhf{}
\fancyhead[L]{\small\color{deitelblue}\textbf{C: Como Programar -- 6ª ed. | Deitel \& Deitel}}
\fancyhead[R]{\small\color{deitelblue}Cap. 2 -- Fazendo a Diferença}
\fancyfoot[C]{\small\thepage}
\renewcommand{\headrulewidth}{0.4pt}

\titleformat{\section}[block]{\large\bfseries\color{deitelblue}}{\thesection.}{0.5em}{}[\titlerule]
\titleformat{\subsection}[block]{\normalsize\bfseries\color{deitelblue!80}}{\thesubsection.}{0.5em}{}

\hypersetup{colorlinks=true,linkcolor=deitelblue,urlcolor=ecogreen}

\begin{document}

\begin{titlepage}
  \centering\vspace*{2cm}
  {\color{deitelblue}\rule{\linewidth}{2pt}}\\[0.4cm]
  {\huge\bfseries\color{deitelblue}C: Como Programar}\\[0.2cm]
  {\large\color{deitelblue}6ª Edição $\cdot$ Paul Deitel \& Harvey Deitel}\\[0.2cm]
  {\color{deitelblue}\rule{\linewidth}{2pt}}\\[1cm]
  {\LARGE\bfseries Capítulo 2}\\[0.3cm]
  {\Large\color{deitelblue!80}Introdução à Programação em C}\\[0.8cm]
  \begin{tcolorbox}[colback=ecolightgreen,colframe=ecogreen,width=0.85\linewidth]
    \centering{\large\bfseries\color{ecogreen}Fazendo a Diferença}\\[0.3cm]
    {\normalsize Exercícios 2.32 (Calculadora de IMC) e 2.33 (Transporte Solidário)\\[0.2cm]
    Programas completos em C com análise didática profunda}
  \end{tcolorbox}
  \vfill{\small\color{deitelblue}``Encorajamos você a usar computadores e a Internet para
  pesquisar e resolver problemas que realmente fazem a diferença.'' -- Deitel \& Deitel}
\end{titlepage}

\tableofcontents\newpage

\section*{Sobre este Capítulo e suas Limitações}

No Capítulo~2, aprendemos a programar com \textbf{variáveis inteiras} (\texttt{int}),
funções \texttt{printf}/\texttt{scanf}, operadores aritméticos e a instrução \texttt{if}.
Ainda não temos acesso a:
\begin{itemize}
  \item \textbf{Ponto flutuante} (\texttt{float}/\texttt{double}) -- será estudado no Cap.~3
  \item \textbf{\texttt{if...else}} -- será estudado no Cap.~3
  \item \textbf{Laços} (\texttt{while}, \texttt{for}) -- serão estudados nos Caps.~3 e~4
\end{itemize}

\begin{atencao}
  Os exercícios ``Fazendo a Diferença'' do Cap.~2 pedem programas que,
  idealmente, usariam \texttt{double}. O próprio enunciado reconhece isso e
  instrui o aluno a usar \texttt{int} por ora, alertando que os resultados
  serão inteiros (sem casas decimais). Apresentamos duas versões:
  a \textbf{versão Cap.~2} (apenas \texttt{int}) e uma
  \textbf{versão aprimorada} (com \texttt{double}, antecipando o Cap.~3)
  para ilustrar a diferença e motivar o aprendizado.
\end{atencao}

\newpage

% ============================================================
\section{Exercício 2.32 -- Calculadora de Índice de Massa Corporal (IMC)}
% ============================================================

\begin{enunciadoBox}[Enunciado 2.32]
\textbf{Calculadora de Índice de Massa Corporal.}
A fórmula para calcular o IMC é:
\[
  IMC = \frac{\text{pesoEmQuilos}}{\text{alturaEmMetros} \times \text{alturaEmMetros}}
\]
Crie uma aplicação que leia o peso (em kg) e a altura (em metros), calcule e apresente
o IMC. Exiba também as faixas do Ministério da Saúde para o usuário avaliar seu IMC.
[\textit{Nota: neste capítulo, use \texttt{int}; resultados serão inteiros.
No Cap.~4, usaremos \texttt{double} para resultados com casas decimais.}]
\end{enunciadoBox}

\bigskip

\subsection{Análise do Problema}

O IMC foi introduzido no Exercício~1.11 (Capítulo~1). Agora, finalmente implementamos
o programa em C. Recapitulando a fórmula:

\[
  IMC = \frac{m}{h^2}
\]

onde $m$ é a massa em quilogramas e $h$ é a altura em metros. O resultado é adimensional
(kg/m²) e classificado segundo as faixas da OMS/Ministério da Saúde.

\subsection{Tabela de Classificação (Ministério da Saúde)}

\begin{center}
\renewcommand{\arraystretch}{1.5}
\begin{tabular}{c l}
  \toprule \textbf{IMC (kg/m²)} & \textbf{Classificação} \\ \midrule
  Menor que 18,5 & Abaixo do peso \\
  Entre 18,5 e 24,9 & Normal (saudável) \\
  Entre 25,0 e 29,9 & Acima do peso \\
  30 ou mais & Obeso \\ \bottomrule
\end{tabular}
\end{center}

\subsection{Versão Capítulo 2 -- com \texttt{int} (conforme enunciado)}

\begin{respostabox}
\begin{lstlisting}
/* Ex2.32 (Versao Cap.2): calculadora_imc_int.c
   Calcula o IMC usando apenas tipo int (inteiros)
   NOTA: resultados serao inteiros -- sem casas decimais
   Para resultados precisos, use double (Cap. 3/4) */
#include <stdio.h>

int main( void )
{
    int peso;     /* peso em quilogramas (inteiro) */
    int altura;   /* altura em metros (inteiro)
                     LIMITACAO: nao aceita ex: 1.75 m */
    int imc;      /* resultado do calculo do IMC (inteiro) */

    /* Solicita dados ao usuario */
    printf( "=== Calculadora de IMC ===\n" );
    printf( "Digite seu peso em kg (numero inteiro): " );
    scanf( "%d", &peso );

    printf( "Digite sua altura em metros (numero inteiro): " );
    scanf( "%d", &altura );

    /* Calcula o IMC: peso / (altura * altura) */
    imc = peso / ( altura * altura );

    /* Exibe o resultado */
    printf( "\nSeu IMC calculado (inteiro): %d kg/m2\n", imc );

    /* Exibe a tabela de referencia */
    printf( "\n----- VALORES DE IMC (Ministerio da Saude) -----\n" );
    printf( "Abaixo do peso:   menor que 18,5\n" );
    printf( "Normal:           entre 18,5 e 24,9\n" );
    printf( "Acima do peso:    entre 25 e 29,9\n" );
    printf( "Obeso:            30 ou mais\n" );
    printf( "------------------------------------------------\n" );

    /* Classificacao com base no valor inteiro do IMC */
    if ( imc < 19 ) {
        printf( "Classificacao: Abaixo do peso\n" );
    }
    if ( imc >= 19 ) {
        if ( imc < 25 ) {
            printf( "Classificacao: Normal\n" );
        }
    }
    if ( imc >= 25 ) {
        if ( imc < 30 ) {
            printf( "Classificacao: Acima do peso\n" );
        }
    }
    if ( imc >= 30 ) {
        printf( "Classificacao: Obeso\n" );
    }

    printf( "\nATENCAO: Consulte sempre um medico para avaliacao.\n" );

    return 0;
} /* fim da funcao main */
\end{lstlisting}

\begin{saida}
=== Calculadora de IMC ===\\
Digite seu peso em kg (numero inteiro): 80\\
Digite sua altura em metros (numero inteiro): 2\\
\\
Seu IMC calculado (inteiro): 20 kg/m2\\
\\
----- VALORES DE IMC (Ministerio da Saude) -----\\
Abaixo do peso:   menor que 18,5\\
Normal:           entre 18,5 e 24,9\\
Acima do peso:    entre 25 e 29,9\\
Obeso:            30 ou mais\\
------------------------------------------------\\
Classificacao: Normal\\
\\
ATENCAO: Consulte sempre um medico para avaliacao.
\end{saida}
\end{respostabox}

\subsection{Limitação da Versão com \texttt{int} e Motivação para o Cap.~3}

\begin{atencao}
A versão com \texttt{int} tem uma limitação séria: não consegue representar alturas
como \texttt{1.75} metros (já que \texttt{int} não aceita decimais). Ao digitar
\texttt{1} como altura, o denominador é \texttt{1 * 1 = 1}, e o IMC seria igual ao
peso -- completamente errado! Isso demonstra por que o Capítulo~3 é fundamental.
\end{atencao}

\subsection{Versão Aprimorada -- com \texttt{double} (antecipando o Capítulo~3)}

\begin{respostabox}
\begin{lstlisting}
/* Ex2.32 (Versao Aprimorada): calculadora_imc_double.c
   Usa tipo double para representar numeros fracionarios
   Conteudo do Cap. 3 -- apresentado aqui como antecipacao */
#include <stdio.h>

int main( void )
{
    double peso;   /* peso em kg -- aceita fracionario */
    double altura; /* altura em metros -- aceita 1.75, 1.80, etc. */
    double imc;    /* IMC calculado com precisao */

    printf( "=== Calculadora de IMC (versao precisa) ===\n" );
    printf( "Digite seu peso em kg (ex: 75.5): " );
    scanf( "%lf", &peso );   /* %lf le um double */

    printf( "Digite sua altura em metros (ex: 1.75): " );
    scanf( "%lf", &altura );

    /* Calcula o IMC com precisao de ponto flutuante */
    imc = peso / ( altura * altura );

    printf( "\nSeu IMC: %.2f kg/m2\n", imc ); /* 2 casas decimais */

    /* Exibe a tabela de referencia */
    printf( "\n----- VALORES DE IMC (Ministerio da Saude) -----\n" );
    printf( "Abaixo do peso:   menor que 18,5\n" );
    printf( "Normal:           entre 18,5 e 24,9\n" );
    printf( "Acima do peso:    entre 25 e 29,9\n" );
    printf( "Obeso:            30 ou mais\n" );
    printf( "------------------------------------------------\n" );

    /* Classificacao precisa */
    if ( imc < 18.5 ) {
        printf( "Classificacao: Abaixo do peso\n" );
    }
    if ( imc >= 18.5 ) {
        if ( imc < 25.0 ) {
            printf( "Classificacao: Normal\n" );
        }
    }
    if ( imc >= 25.0 ) {
        if ( imc < 30.0 ) {
            printf( "Classificacao: Acima do peso\n" );
        }
    }
    if ( imc >= 30.0 ) {
        printf( "Classificacao: Obeso\n" );
    }

    printf( "\nATENCAO: O IMC e uma medida de triagem.\n" );
    printf( "Consulte sempre um profissional de saude.\n" );

    return 0;
} /* fim da funcao main */
\end{lstlisting}

\begin{saida}
=== Calculadora de IMC (versao precisa) ===\\
Digite seu peso em kg (ex: 75.5): 75.0\\
Digite sua altura em metros (ex: 1.75): 1.75\\
\\
Seu IMC: 24.49 kg/m2\\
\\
----- VALORES DE IMC (Ministerio da Saude) -----\\
Abaixo do peso:   menor que 18,5\\
Normal:           entre 18,5 e 24,9\\
Acima do peso:    entre 25 e 29,9\\
Obeso:            30 ou mais\\
------------------------------------------------\\
Classificacao: Normal\\
\\
ATENCAO: O IMC e uma medida de triagem.\\
Consulte sempre um profissional de saude.
\end{saida}
\end{respostabox}

\subsection{Conexão com o Capítulo~1 e Progressão do Livro}

Este exercício completa o ciclo iniciado no Exercício~1.11 (Capítulo~1), onde
exploramos as calculadoras de IMC on-line e estudamos suas fórmulas. Agora implementamos
o programa em C. A progressão ilustra perfeitamente a abordagem incremental do livro:
\begin{enumerate}
  \item \textbf{Cap.~1:} Entendemos o problema, pesquisamos as fórmulas, desenvolvemos o algoritmo
  \item \textbf{Cap.~2:} Implementamos o programa com as ferramentas disponíveis (\texttt{int}, \texttt{if})
  \item \textbf{Cap.~3:} Usaremos \texttt{double} e \texttt{if...else} para uma versão mais elegante
  \item \textbf{Cap.~4:} Adicionaremos \texttt{while} para validar entradas inválidas
\end{enumerate}

\newpage

% ============================================================
\section{Exercício 2.33 -- Calculadora de Economias com Transporte Solidário}
% ============================================================

\begin{enunciadoBox}[Enunciado 2.33]
\textbf{Calculadora de economias com o transporte solidário.}
Pesquise websites sobre transporte solidário (\textit{carpool}). Crie uma aplicação
que calcule a despesa diária com o automóvel, para que o usuário estime quanto
poderia economizar com o transporte solidário. A aplicação deve solicitar:
\begin{enumerate}[label=\alph*)]
  \item Total de quilômetros dirigidos por dia
  \item Custo por litro de combustível
  \item Média de quilômetros por litro
  \item Preço de estacionamento por dia
  \item Gastos diários com pedágios
\end{enumerate}
\end{enunciadoBox}

\subsection{Análise do Problema e Pesquisa}

\begin{pesquisa}
Este exercício requer pesquisa sobre transporte solidário. O \textit{carpool}
(ou transporte solidário compartilhado) é uma prática em que duas ou mais pessoas
compartilham um único veículo para realizar o mesmo trajeto. Além da economia financeira,
reduz emissões de CO\textsubscript{2} e congestionamentos.

\textbf{Fórmula do custo diário com combustível:}
\[
  \text{custo\_combustivel} = \frac{\text{km\_por\_dia}}{\text{km\_por\_litro}} \times \text{preco\_litro}
\]

\textbf{Custo total diário:}
\[
  \text{custo\_total} = \text{custo\_combustivel} + \text{estacionamento} + \text{pedagios}
\]

\textbf{Economia com carpool (2 pessoas):}
\[
  \text{economia} = \frac{\text{custo\_total}}{2}
\]
\end{pesquisa}

\subsection{Versão Capítulo 2 -- com \texttt{int} (valores inteiros)}

\begin{respostabox}
\begin{lstlisting}
/* Ex2.33 (Versao Cap.2): transporte_solidario_int.c
   Calcula despesa diaria com automovel usando apenas int
   NOTA: resultados sem casas decimais */
#include <stdio.h>

int main( void )
{
    int kmPorDia;          /* quilometros dirigidos por dia */
    int precoPorLitro;     /* preco do combustivel em centavos */
    int kmPorLitro;        /* consumo medio do veiculo */
    int estacionamento;    /* preco do estacionamento em centavos */
    int pedagios;          /* total de pedagios em centavos */

    int litrosPorDia;      /* combustivel consumido diariamente */
    int custoCombustivel;  /* custo total do combustivel */
    int custoTotal;        /* custo total diario */
    int economiaCarpool;   /* economia com carpool (2 pessoas) */

    printf( "=== Calculadora de Transporte Solidario ===\n\n" );
    printf( "Informe seus dados diarios de conducao:\n" );

    printf( "Total de km dirigidos por dia: " );
    scanf( "%d", &kmPorDia );

    printf( "Preco do combustivel (centavos por litro, ex: 590 = R$5,90): " );
    scanf( "%d", &precoPorLitro );

    printf( "Media de km por litro do seu veiculo: " );
    scanf( "%d", &kmPorLitro );

    printf( "Preco do estacionamento por dia (centavos): " );
    scanf( "%d", &estacionamento );

    printf( "Gastos com pedagios por dia (centavos): " );
    scanf( "%d", &pedagios );

    /* Calcula o consumo de combustivel */
    litrosPorDia    = kmPorDia / kmPorLitro;
    custoCombustivel = litrosPorDia * precoPorLitro;
    custoTotal      = custoCombustivel + estacionamento + pedagios;
    economiaCarpool = custoTotal / 2;  /* divide com 1 colega */

    /* Exibe os resultados */
    printf( "\n=== CUSTO DIARIO COM AUTOMOVEL ===\n" );
    printf( "Combustivel:     %d centavos\n", custoCombustivel );
    printf( "Estacionamento:  %d centavos\n", estacionamento );
    printf( "Pedagios:        %d centavos\n", pedagios );
    printf( "Total diario:    %d centavos\n", custoTotal );
    printf( "-----------------------------------\n" );
    printf( "Com carpool (2 pessoas), voce pagaria: %d centavos/dia\n",
            economiaCarpool );
    printf( "Economia diaria:   %d centavos\n", custoTotal - economiaCarpool );
    printf( "Economia mensal:   %d centavos (22 dias uteis)\n",
            ( custoTotal - economiaCarpool ) * 22 );

    return 0;
} /* fim da funcao main */
\end{lstlisting}

\begin{saida}
=== Calculadora de Transporte Solidario ===\\
\\
Informe seus dados diarios de conducao:\\
Total de km dirigidos por dia: 40\\
Preco do combustivel (centavos por litro, ex: 590 = R\$5,90): 590\\
Media de km por litro do seu veiculo: 10\\
Preco do estacionamento por dia (centavos): 2000\\
Gastos com pedagios por dia (centavos): 500\\
\\
=== CUSTO DIARIO COM AUTOMOVEL ===\\
Combustivel:     2360 centavos\\
Estacionamento:  2000 centavos\\
Pedagios:        500 centavos\\
Total diario:    4860 centavos\\
-----------------------------------\\
Com carpool (2 pessoas), voce pagaria: 2430 centavos/dia\\
Economia diaria:   2430 centavos\\
Economia mensal:   53460 centavos (22 dias uteis)
\end{saida}
\end{respostabox}

\subsection{Versão Aprimorada -- com \texttt{double} e valores reais}

\begin{respostabox}
\begin{lstlisting}
/* Ex2.33 (Versao Aprimorada): transporte_solidario_double.c
   Usa double para calculos precisos com valores reais
   Antecipacao do Cap. 3 */
#include <stdio.h>

int main( void )
{
    double kmPorDia;          /* quilometros por dia */
    double precoPorLitro;     /* preco do litro de combustivel (R$) */
    double kmPorLitro;        /* eficiencia do veiculo */
    double estacionamento;    /* custo do estacionamento (R$) */
    double pedagios;          /* custo com pedagios (R$) */

    double custoCombustivel;  /* custo diario com combustivel */
    double custoTotal;        /* custo diario total */
    double custoCarpool2;     /* custo com 2 pessoas */
    double custoCarpool3;     /* custo com 3 pessoas */
    double custoCarpool4;     /* custo com 4 pessoas */

    printf( "=== Calculadora de Transporte Solidario (Versao Precisa) ===\n\n" );
    printf( "Total de km dirigidos por dia: " );
    scanf( "%lf", &kmPorDia );

    printf( "Preco do combustivel por litro (R$): " );
    scanf( "%lf", &precoPorLitro );

    printf( "Media de km por litro do seu veiculo: " );
    scanf( "%lf", &kmPorLitro );

    printf( "Preco do estacionamento por dia (R$): " );
    scanf( "%lf", &estacionamento );

    printf( "Gastos com pedagios por dia (R$): " );
    scanf( "%lf", &pedagios );

    /* Calculos precisos */
    custoCombustivel = ( kmPorDia / kmPorLitro ) * precoPorLitro;
    custoTotal       = custoCombustivel + estacionamento + pedagios;
    custoCarpool2    = custoTotal / 2.0;
    custoCarpool3    = custoTotal / 3.0;
    custoCarpool4    = custoTotal / 4.0;

    /* Relatorio detalhado */
    printf( "\n======================================\n" );
    printf( "   RELATORIO DE CUSTOS DIARIOS\n" );
    printf( "======================================\n" );
    printf( "Combustivel:       R$ %.2f\n", custoCombustivel );
    printf( "Estacionamento:    R$ %.2f\n", estacionamento );
    printf( "Pedagios:          R$ %.2f\n", pedagios );
    printf( "--------------------------------------\n" );
    printf( "Custo total (solo): R$ %.2f\n", custoTotal );
    printf( "======================================\n" );
    printf( "\n   ECONOMIA COM TRANSPORTE SOLIDARIO\n" );
    printf( "--------------------------------------\n" );
    printf( "Com 2 pessoas:  R$ %.2f/pessoa/dia\n", custoCarpool2 );
    printf( "  Voce economiza: R$ %.2f/dia\n", custoTotal - custoCarpool2 );
    printf( "With 3 pessoas: R$ %.2f/pessoa/dia\n", custoCarpool3 );
    printf( "  Voce economiza: R$ %.2f/dia\n", custoTotal - custoCarpool3 );
    printf( "Com 4 pessoas:  R$ %.2f/pessoa/dia\n", custoCarpool4 );
    printf( "  Voce economiza: R$ %.2f/dia\n", custoTotal - custoCarpool4 );
    printf( "--------------------------------------\n" );
    printf( "Economia ANUAL com 2 pessoas (250 dias): R$ %.2f\n",
            ( custoTotal - custoCarpool2 ) * 250.0 );
    printf( "======================================\n" );

    return 0;
} /* fim da funcao main */
\end{lstlisting}

\begin{saida}
=== Calculadora de Transporte Solidario (Versao Precisa) ===\\
\\
Total de km dirigidos por dia: 40.0\\
Preco do combustivel por litro (R\$): 5.90\\
Media de km por litro do seu veiculo: 10.5\\
Preco do estacionamento por dia (R\$): 20.00\\
Gastos com pedagios por dia (R\$): 5.00\\
\\
======================================\\
   RELATORIO DE CUSTOS DIARIOS\\
======================================\\
Combustivel:       R\$ 22.48\\
Estacionamento:    R\$ 20.00\\
Pedagios:          R\$ 5.00\\
--------------------------------------\\
Custo total (solo): R\$ 47.48\\
======================================\\
\\
   ECONOMIA COM TRANSPORTE SOLIDARIO\\
--------------------------------------\\
Com 2 pessoas:  R\$ 23.74/pessoa/dia\\
  Voce economiza: R\$ 23.74/dia\\
Com 3 pessoas: R\$ 15.83/pessoa/dia\\
  Voce economiza: R\$ 31.65/dia\\
Com 4 pessoas:  R\$ 11.87/pessoa/dia\\
  Voce economiza: R\$ 35.61/dia\\
--------------------------------------\\
Economia ANUAL com 2 pessoas (250 dias): R\$ 5935.00\\
======================================
\end{saida}
\end{respostabox}

\subsection{Impacto Ambiental -- Conexão com o Exercício 1.10}

\begin{tcolorbox}[colback=ecolightgreen,colframe=ecogreen,
  title={\bfseries Benefícios Ambientais do Transporte Solidário}]

O Exercício~1.10 (Calculadora de Emissão de Carbono) nos mostrou que o
transporte por carro a gasolina emite aproximadamente 2,31 kg de CO\textsubscript{2}
por litro. No exemplo acima, com 40 km/dia a 10,5 km/litro, o usuário consome
$\approx 3,81$ litros/dia, emitindo $\approx 8,8$ kg de CO\textsubscript{2}/dia.

Com 2 pessoas no mesmo carro:
\begin{itemize}
  \item As emissões \textit{por pessoa} caem para $\approx 4,4$ kg CO\textsubscript{2}/dia
  \item Em 250 dias de trabalho: redução de $\approx 1.100$ kg CO\textsubscript{2}/ano por pessoa
  \item Equivalente a plantar $\approx 50$ árvores adultas por ano
\end{itemize}

Os exercícios ``Fazendo a Diferença'' do Capítulo~2 demonstram que programação
em C não é apenas técnica -- é uma ferramenta poderosa para entender e impactar o mundo.
\end{tcolorbox}

\newpage

\section*{Reflexão -- Progressão Incremental Deitel}

\begin{tcolorbox}[colback=deitellightblue,colframe=deitelblue,
  title={\bfseries O que aprendemos e o que está por vir}]

\textbf{Neste Capítulo~2:}
\begin{itemize}[noitemsep]
  \item Escrevemos nossos primeiros programas reais em C
  \item Usamos \texttt{printf}/\texttt{scanf} para comunicação com o usuário
  \item Trabalhamos com variáveis \texttt{int} e operadores aritméticos
  \item Tomamos decisões simples com \texttt{if}
  \item Compreendemos os conceitos de memória destrutiva e não destrutiva
\end{itemize}

\medskip\textbf{No Capítulo~3 (próximo passo), aprenderemos:}
\begin{itemize}[noitemsep]
  \item Números de ponto flutuante (\texttt{float}/\texttt{double})
  \item \texttt{if...else} para decisões binárias elegantes
  \item A estrutura de repetição \texttt{while}
  \item Algoritmos com refinamentos sucessivos (\textit{top-down})
\end{itemize}

\medskip\textbf{Com essas ferramentas, as calculadoras de IMC e transporte solidário
ficarão muito mais precisas e robustas.}

\begin{boaprática}
  A abordagem incremental do livro é intencional: você aprende um conjunto
  de ferramentas, as aplica em problemas reais, e então recebe novas ferramentas
  para resolver problemas mais sofisticados. Cada novo capítulo expande sua
  caixa de ferramentas -- e sua capacidade de fazer a diferença.
\end{boaprática}

\end{tcolorbox}

\end{document}
